\chapter{Methodology}
\label{chap:methodology}

\chapterabstract{This chapter describes the proposed zero-shot static human-scene interaction synthesis pipeline. The method converts a natural-language instruction, a calibrated ScanNet++ RGB view, and a reconstructed scene mesh into a single SMPL-X body placed in the original scene coordinate frame.}

\newcommand{\methodbox}[2]{%
  \fcolorbox{ThesisGray}{ThesisLightGray}{%
    \begin{minipage}[c][2.0cm][c]{#1}
      \centering\small #2
    \end{minipage}%
  }%
}

\section{Overview}
\label{sec:method-overview}

This chapter presents the proposed zero-shot pipeline for static 3D human--scene interaction synthesis. The method takes as input a natural-language interaction instruction $y$, a calibrated RGB view $I$ from a ScanNet++ scene, the corresponding camera intrinsics $K$, camera rotation $R$, camera translation $t$, and the reconstructed scene mesh $M_s$. The output is a single SMPL-X body placed in the ScanNet++ scene coordinate frame. The synthesized result is static: it represents the interaction state, or a contact-ready pre-action state, rather than a full motion sequence or an object-state change.

The method is designed to handle both underspecified and detailed interaction instructions. For example, the instruction ``a person lifting a table'' specifies the interaction at a high level, leaving the exact body configuration and contact layout to be inferred. Functional interactions are more demanding: an instruction such as ``a person opening a washing machine'' requires the method to infer that contact should occur near the functional door or handle region, not merely somewhere on the washing machine object. The method also supports more detailed instructions, such as ``a person on the treadmill with both hands on the side frames and one foot on the belt,'' where the user explicitly constrains the pose and the required body-part contacts. In all cases, the central challenge is the same: the language instruction must be converted into metric 3D contact constraints that are consistent with the reconstructed scene.

A visually plausible generated image is not sufficient for this task, because the final body must also be geometrically compatible with the scene mesh. A generated person may appear reasonable in the input view while still floating above the floor, intersecting furniture, or contacting the wrong object part in 3D. The proposed method therefore combines foundation-model reasoning with explicit geometric grounding. Foundation models are used to infer semantic and visual interaction cues, while camera geometry and optimization convert these cues into explicit 3D contact targets and a physically plausible SMPL-X body.

\Cref{fig:method-pipeline} illustrates the complete pipeline. The method consists of five main modules. First, the instruction and scene image are converted into a Scene Interaction Graph (SIG), which identifies the target object or floor, the relevant human body parts, and the required body-part--scene contact edges. Second, an image generation model inserts a plausible human into the original scene view according to the interaction instruction. This generated human frame acts as a visual hypothesis for body pose, orientation, laterality, and approximate contact layout. Third, the generated human frame, the original scene image, and the SIG are used to estimate scene-side contact regions. Each required contact region is marked on the original scene image using a distinct bright color, producing body-part-specific contact masks. When the interaction involves floor support, floor regions are obtained separately using SAM3. The accepted contact masks are then projected to the reconstructed scene mesh to form 3D scene contact targets.

In parallel, the fourth module initializes the human body. The generated human frame is passed to GVHMR, which predicts an initial SMPL-X body. This initialization provides a plausible articulated pose and global body configuration, but it is not assumed to satisfy the scene contact constraints exactly. The fifth module refines the initialized body through scene-aware optimization. During this stage, the scene geometry and projected contact targets are kept fixed, while the SMPL-X translation, global orientation, and body pose are optimized. The objective combines contact attraction, penetration avoidance, and priors that keep the result close to the GVHMR initialization.

The key idea of the proposed approach is that image generation is used not only to synthesize a plausible human interaction image, but also to localize fine-grained scene-side contact regions. These contact annotations are not treated as final outputs. Instead, they are converted into explicit 3D contact targets on the reconstructed mesh and used to guide SMPL-X optimization. In this way, the method combines the flexibility of open-vocabulary foundation models with the geometric accountability required for static 3D human--scene interaction synthesis.

The full prompts used for the foundation-model stages are provided in \cref{app:system-prompts} for reproducibility. The methodology chapter describes the algorithmic role of each model call, while the appendix records the exact prompt wording used by the implementation.

\begin{figure}[t]
  \centering

  % TODO: Replace this placeholder with the final methodology diagram.
  % Recommended layout: one horizontal pipeline using a single running example.
  %
  % Suggested panels:
  % (a) Input ScanNet++ RGB view + instruction
  % (b) Scene Interaction Graph (SIG)
  % (c) Human-inpainted frame
  % (d) Colored contact annotations on the original scene image
  % (e) Projected 3D contact targets on the scene mesh
  % (f) GVHMR SMPL-X initialization
  % (g) Optimized SMPL-X body in the ScanNet++ scene
  %
  % Suggested examples:
  % - Functional: ``a person opening a washing machine''
  % - Detailed: ``a person on the treadmill with both hands on the side frames and one foot on the belt''

  \fbox{
    \begin{minipage}[c][0.28\textheight][c]{0.96\textwidth}
      \centering
      \textbf{Methodology Pipeline Placeholder}\\[1.0em]

      \begin{tabular}{c@{\quad$\rightarrow$\quad}c@{\quad$\rightarrow$\quad}c@{\quad$\rightarrow$\quad}c}
        \begin{minipage}{0.18\textwidth}
          \centering
          Input\\
          \small RGB view $I$\\
          instruction $y$\\
          mesh $M_s$
        \end{minipage}
        &
        \begin{minipage}{0.18\textwidth}
          \centering
          SIG\\
          \small body parts\\
          object/floor\\
          contact edges
        \end{minipage}
        &
        \begin{minipage}{0.18\textwidth}
          \centering
          Human frame\\
          \small inpainted\\
          interaction\\
          image
        \end{minipage}
        &
        \begin{minipage}{0.18\textwidth}
          \centering
          Contact masks\\
          \small colored regions\\
          on original\\
          scene image
        \end{minipage}
      \end{tabular}

      \vspace{1.2em}

      \begin{tabular}{c@{\quad$\rightarrow$\quad}c@{\quad$\rightarrow$\quad}c}
        \begin{minipage}{0.22\textwidth}
          \centering
          3D contact targets\\
          \small projected to\\
          ScanNet++ mesh
        \end{minipage}
        &
        \begin{minipage}{0.22\textwidth}
          \centering
          GVHMR initialization\\
          \small initial\\
          SMPL-X body
        \end{minipage}
        &
        \begin{minipage}{0.22\textwidth}
          \centering
          Optimization\\
          \small final SMPL-X\\
          in scene coordinates
        \end{minipage}
      \end{tabular}
    \end{minipage}
  }

  \caption{
  Overview of the proposed zero-shot static HSI synthesis pipeline. Given a natural-language instruction, a calibrated ScanNet++ RGB view, and the reconstructed scene mesh, the method generates a Scene Interaction Graph, synthesizes a human interaction image, estimates body-part-specific scene contact regions on the original scene image, projects these contacts to the scene mesh, initializes an SMPL-X body with GVHMR, and refines the body through scene-aware optimization.
  }
  \label{fig:method-pipeline}
\end{figure}


\section{Problem Formulation}
\label{sec:problem-formulation}

Given a natural-language interaction instruction, a calibrated RGB image, and a reconstructed 3D scene, the goal is to synthesize a single human body that realizes the requested interaction in the original scene coordinate frame. Let $I$ denote the selected RGB view from a ScanNet++ scene. The image has spatial resolution
\begin{equation}
\label{eq:input-image}
I \in \mathbb{R}^{H \times W \times 3}.
\end{equation}

The camera calibration is defined by the intrinsic matrix $K$, the camera rotation $R$, and the camera translation $t$. The reconstructed scene mesh is denoted by
\begin{equation}
\label{eq:scene-mesh}
M_s = (V_s, F_s),
\end{equation}
where $V_s$ are the scene vertices in the ScanNet++ world coordinate frame and $F_s$ are the triangular faces. The scene vertices satisfy
\begin{equation}
\label{eq:scene-vertices}
V_s \in \mathbb{R}^{N_s \times 3}.
\end{equation}

The user instruction is denoted by $y$. It can be underspecified, such as ``a person lifting a table,'' functional, such as ``a person opening a washing machine,'' or detailed, such as ``a person on the treadmill with both hands on the side frames and one foot on the belt.'' In all cases, the method must infer a body configuration that is consistent with both the language instruction and the reconstructed scene geometry.

The output of the method is a posed SMPL-X body in the same world coordinate frame as the reconstructed scene. SMPL-X represents the body as a differentiable parametric mesh. We denote the body parameters by
\begin{equation}
\label{eq:smplx-params}
\theta =
(
\mathbf{T},
\mathbf{r},
\mathbf{p},
\boldsymbol{\beta}
).
\end{equation}
Here, $\mathbf{T}$ is the global translation, $\mathbf{r}$ is the global orientation, $\mathbf{p}$ is the articulated body pose, and $\boldsymbol{\beta}$ are the body shape parameters. Applying the SMPL-X model to these parameters produces the human mesh
\begin{equation}
\label{eq:human-mesh}
M_h(\theta) = (V_h(\theta), F_h),
\end{equation}
where $V_h(\theta)$ are the posed body vertices and $F_h$ are the fixed SMPL-X mesh faces. The posed body vertices satisfy
\begin{equation}
\label{eq:human-vertices}
V_h(\theta) \in \mathbb{R}^{N_h \times 3}.
\end{equation}

The synthesis target is not a motion sequence. The method estimates one static interaction state, or a contact-ready state, that expresses the instruction. For object-manipulation prompts such as opening a washing machine or lifting a table, the object itself remains fixed in the reconstructed scene. The goal is to synthesize the body configuration that makes the interaction visually and geometrically plausible, not to simulate subsequent object motion or state change.

The instruction $y$ is converted into a set of body--scene contact requirements. These requirements are represented by a Scene Interaction Graph,
\begin{equation}
\label{eq:sig}
G = (\mathcal{H}, \mathcal{S}, \mathcal{E}),
\end{equation}
where $\mathcal{H}$ is the set of human body-part nodes, $\mathcal{S}$ is the set of scene nodes, and $\mathcal{E}$ is the set of contact or support edges. Human nodes correspond to supported body parts such as hands, arms, feet, legs, hips, back, and head. Scene nodes correspond either to the target object or to the floor. Each edge is written as
\begin{equation}
\label{eq:sig-edge}
e = (h_e, s_e),
\qquad
e \in \mathcal{E},
\end{equation}
and indicates that body part $h_e$ should contact scene element $s_e$.

Each human node $h_e$ is associated with a fixed subset of SMPL-X contact vertices. We denote this body-side contact set by
\begin{equation}
\label{eq:body-contact-set}
B_e \subset \{1,\ldots,N_h\}.
\end{equation}
For a given body configuration $\theta$, the corresponding body-side contact points are
\begin{equation}
\label{eq:body-contact-points}
P_e^h(\theta)
=
\{V_h(\theta)_i \mid i \in B_e\}.
\end{equation}

The scene-side contact target is not manually annotated. Instead, it is estimated from generated visual contact annotations and projected to the reconstructed scene mesh. For each edge $e$, this produces a set of scene contact vertices or sampled surface points, denoted by
\begin{equation}
\label{eq:scene-contact-target}
P_e^s \subset \mathbb{R}^{3}.
\end{equation}

The optimization problem can therefore be stated as finding SMPL-X parameters $\theta$ such that the body-side contact points $P_e^h(\theta)$ are close to the corresponding scene-side contact targets $P_e^s$, while the body remains physically plausible and avoids severe penetration with the scene. The body is initialized using GVHMR from the generated human frame. Let the initial SMPL-X estimate be
\begin{equation}
\label{eq:smplx-init-params}
\theta_0 =
(
\mathbf{T}_0,
\mathbf{r}_0,
\mathbf{p}_0,
\boldsymbol{\beta}_0
).
\end{equation}

During optimization, only the global translation $\mathbf{T}$, global orientation $\mathbf{r}$, and body pose $\mathbf{p}$ are updated. The shape parameters $\boldsymbol{\beta}$ are initialized from GVHMR and kept fixed.

The final objective has the following general form:
\begin{equation}
\label{eq:method-total-loss}
\mathcal{L}_{\mathrm{total}}(\theta)
=
\lambda_{\mathrm{contact}} \mathcal{L}_{\mathrm{contact}}(\theta)
+
\lambda_{\mathrm{penetration}} \mathcal{L}_{\mathrm{penetration}}(\theta)
+
\lambda_{\mathrm{pose}} \mathcal{L}_{\mathrm{pose}}(\theta)
+
\lambda_{\mathrm{orient}} \mathcal{L}_{\mathrm{orient}}(\theta).
\end{equation}

The final body parameters are obtained by minimizing this objective with respect to the global translation, global orientation, and body pose:
\begin{equation}
\label{eq:method-optimization-problem}
\theta^{\ast}
=
\operatorname*{arg\,min}_{\mathbf{T},\,\mathbf{r},\,\mathbf{p}}
\mathcal{L}_{\mathrm{total}}(\theta).
\end{equation}

The contact term pulls the specified SMPL-X body regions toward the projected scene contact targets. The penetration term discourages intersections between the human body and the reconstructed scene. The pose and orientation priors keep the optimized body close to the GVHMR initialization so that contact refinement does not produce unnatural body configurations. The following sections describe how the graph $G$, the human initialization $\theta_0$, and the scene contact targets $P_e^s$ are obtained.

\begin{table}[t]
  \centering
  \caption{Notation used in the methodology chapter.}
  \label{tab:method-notation}
  \begin{tabularx}{\textwidth}{@{}>{\bfseries}p{0.18\textwidth}X@{}}
    \toprule
    Symbol & Meaning \\
    \midrule
    $I$ & Calibrated RGB view from a ScanNet++ scene. \\
    $K,R,t$ & Camera intrinsics, rotation, and translation. \\
    $M_s=(V_s,F_s)$ & Reconstructed scene mesh in the ScanNet++ world coordinate frame. \\
    $y$ & Natural-language interaction instruction. \\
    $G=(\mathcal{H},\mathcal{S},\mathcal{E})$ & Scene Interaction Graph. \\
    $e=(h_e,s_e)$ & Contact edge between body part $h_e$ and scene element $s_e$. \\
    $B_e$ & SMPL-X vertex subset assigned to the body part in edge $e$. \\
    $P_e^h(\theta)$ & Body-side contact points for edge $e$. \\
    $P_e^s$ & Scene-side contact targets projected to the reconstructed scene mesh. \\
    $\theta$ & SMPL-X parameters. \\
    $\theta_0$ & GVHMR initialization. \\
    \bottomrule
  \end{tabularx}
\end{table}


\section{Scene Interaction Graph Generation}
\label{sec:method-sig}

The first module converts the open-language instruction into an explicit set of body--scene contact requirements. The purpose of this module is not to solve the full geometry of the interaction, but to determine which human body parts should contact which scene elements. This intermediate representation is called the Scene Interaction Graph (SIG). It provides the symbolic structure that later guides both contact-region estimation and SMPL-X optimization. The exact system prompt used for SIG generation is provided in \cref{app:system-prompt-sig}; the main text describes only the functional role of the prompt and the validated graph structure.

The SIG is generated from the user instruction $y$ and the input scene image $I$. A language or vision-language model is prompted to identify the target object, the human body parts that are physically important for the interaction, and any required floor support. The output is parsed into a graph
\begin{equation}
\label{eq:sig-generation-output}
G = (\mathcal{H}, \mathcal{S}, \mathcal{E}),
\end{equation}
where $\mathcal{H}$ contains human body-part nodes, $\mathcal{S}$ contains scene nodes, and $\mathcal{E}$ contains contact or support edges. The scene nodes are deliberately restricted to the target object and, when required, the floor. The graph does not require explicit object-part labels such as handle, knob, rung, or edge. Instead, it specifies that a particular body part should contact the target object, while the exact functional contact region is estimated later from image-space contact annotations.

The supported human body-part vocabulary is fixed so that each graph node can be mapped to a predefined SMPL-X contact region. The supported parts are left hand, right hand, left arm, right arm, left leg, right leg, left foot, right foot, head, hips, and back. This restricted vocabulary prevents the graph from introducing body parts that cannot be grounded to SMPL-X vertices. It also keeps the downstream optimization well defined: every edge in the graph corresponds to one known body-side vertex set and one scene-side target.

The graph edges define the required contacts. For example, an underspecified prompt such as ``a person lifting a table'' may produce hand--table edges and foot--floor support edges. A functional prompt such as ``a person opening a washing machine'' may produce a hand--washing-machine edge together with foot--floor support. A detailed prompt such as ``a person on the treadmill with both hands on the side frames and one foot on the belt'' may produce separate left-hand--treadmill, right-hand--treadmill, left-foot--treadmill, and right-foot--floor edges, depending on the specified pose. These examples illustrate that the graph can represent both inferred contacts and user-specified contacts.

\begin{table}[t]
  \centering
  \caption{Example Scene Interaction Graphs for prompts with different levels of specificity.}
  \label{tab:sig-examples}
  \begin{tabularx}{\textwidth}{@{}p{0.29\textwidth}p{0.21\textwidth}X@{}}
    \toprule
    Prompt & Scene nodes & Contact/support edges \\
    \midrule
    ``A person lifting a table.'' &
    target object: table; floor &
    left hand $\rightarrow$ table;
    right hand $\rightarrow$ table;
    left foot $\rightarrow$ floor;
    right foot $\rightarrow$ floor. \\
    \addlinespace
    ``A person opening a washing machine.'' &
    target object: washing machine; floor &
    right hand $\rightarrow$ washing machine;
    left foot $\rightarrow$ floor;
    right foot $\rightarrow$ floor. \\
    \addlinespace
    ``A person on the treadmill with both hands on the side frames and one foot on the belt.'' &
    target object: treadmill; floor &
    left hand $\rightarrow$ treadmill;
    right hand $\rightarrow$ treadmill;
    left foot $\rightarrow$ treadmill;
    right foot $\rightarrow$ floor. \\
    \bottomrule
  \end{tabularx}
\end{table}

After generation, the SIG is validated before being used by later modules. The validation checks that the graph contains at least one target object or floor node, that each human node belongs to the supported body-part vocabulary, and that every edge references a defined scene node. Duplicate edges and unsupported body parts are removed or rejected. If no usable contact edge remains, the pipeline cannot construct meaningful contact targets and the graph must be regenerated.

Each valid graph edge $e = (h_e, s_e)$ is associated with a body-side contact region on SMPL-X. For a human body part $h_e$, the corresponding contact region is a fixed vertex subset $B_e$. These vertex subsets are manually defined from the SMPL-X mesh and represent the surfaces that are most likely to participate in physical contact. For example, hand contacts use palm-side vertices, foot contacts use support-facing foot vertices, and sitting or leaning contacts use appropriate regions on the hips, back, head, or legs. These details are fixed once for the whole method and are reused for all interactions.

The SIG therefore connects language to geometry in a controlled way. It does not determine the final pose, and it does not localize the exact object surface that should be touched. Instead, it specifies the required body-part--scene relations. The following modules use this graph to generate a visual human interaction hypothesis, estimate the corresponding scene-side contact regions, and optimize the SMPL-X body so that the specified body regions contact the projected scene targets.


\section{Human Frame Generation}
\label{sec:method-human-frame}

The second module generates a visual hypothesis of the requested interaction in the input scene view. Given the original ScanNet++ image $I$, the user instruction $y$, and the Scene Interaction Graph $G$, an image generation model is prompted to insert a single human into the scene. The output is an edited image in which the generated person appears to perform the requested interaction while the scene remains otherwise unchanged. The full human-frame generation prompt is reported in \cref{app:system-prompt-human-inpaint}.

The purpose of this module is not to produce the final result. The generated image is used as an intermediate visual prior for two downstream tasks. First, it provides the image evidence needed by GVHMR to initialize an SMPL-X body. Second, it shows the approximate pose, orientation, laterality, and contact layout that the contact-estimation module uses to infer scene-side contact regions. The final output of the method is produced only after contact grounding and 3D SMPL-X optimization.

The image generation prompt is constructed to preserve the original scene as much as possible. The model is instructed to keep the camera viewpoint, room geometry, object layout, lighting, textures, and target object fixed. The prompt also specifies that exactly one person should be inserted. For object-manipulation instructions, the target object should not be moved, opened, lifted, rotated, or deformed by the image model. For example, for an instruction such as ``a person opening a washing machine,'' the desired output is an image of a person reaching for or contacting the washing-machine door or handle region, while the washing machine itself remains in its original closed state. This is important because the reconstructed scene mesh represents the original object geometry and cannot reflect image-only object edits.

The generated human frame must also respect the body-part requirements encoded in the SIG. If the graph contains hand--object edges, the generated person should place the corresponding hands near the target object. If the graph contains foot--floor edges, the person should appear physically supported by the floor. If the instruction explicitly specifies asymmetric contacts, such as one foot on a treadmill belt and the other foot on the floor, the generated frame should preserve this laterality. These requirements make the generated image a useful visual bridge between language-level interaction intent and later metric optimization.

In practice, the image generation model may produce multiple candidate frames. Let these candidates be denoted by
\begin{equation}
\label{eq:human-frame-candidates}
\mathcal{I}_{H}
=
\{
\hat{I}^{(1)},
\hat{I}^{(2)},
\ldots,
\hat{I}^{(N)}
\}.
\end{equation}
A vision-language evaluator is then used to choose the most suitable candidate from this set. The evaluator compares each candidate against the original image, the interaction instruction, and the expected contact structure. It favors candidates that express the requested interaction clearly, preserve the target object and surrounding scene, contain a plausible full-body pose, and expose the relevant body-part contacts for later reasoning. The prompt used for this candidate-selection step is provided in \cref{app:system-prompt-best-frame-selection}.

The selected human frame is denoted by
\begin{equation}
\label{eq:selected-human-frame}
\hat{I}
=
\hat{I}^{(j^\star)}.
\end{equation}
Minor visual artifacts in the selected frame are acceptable because the generated image is not used as a final rendering. However, candidates with severe object distortions, missing body parts, implausible body scale, incorrect laterality, or unrelated interaction geometry are rejected by the evaluator. This selection step improves the stability of the downstream stages because both GVHMR initialization and contact estimation depend on the generated human being visually consistent with the requested interaction.

The selected human frame $\hat{I}$ remains in the same image coordinate system as the original view $I$, which allows it to be compared directly with the original scene image in the contact-estimation module. The generated human in $\hat{I}$ is used only as evidence for pose and contact reasoning. The authoritative scene geometry remains the reconstructed ScanNet++ mesh $M_s$, and the authoritative scene appearance for contact annotation remains the original image $I$.

\begin{figure}[t]
  \centering
  % TODO: Replace with side-by-side images from one example.
  % Suggested panels:
  % (a) original ScanNet++ RGB image
  % (b) generated candidate frames
  % (c) VLM-selected human frame
  % Use either a functional example such as "opening a washing machine"
  % or a detailed pose example such as the treadmill prompt.

  \fbox{
    \begin{minipage}[c][0.22\textheight][c]{0.92\textwidth}
      \centering
      \textbf{Human frame generation placeholder}\\[0.8em]
      Original scene image $I$
      \quad $\rightarrow$ \quad
      candidate human frames $\mathcal{I}_{H}$
      \quad $\rightarrow$ \quad
      selected frame $\hat{I}$
    \end{minipage}
  }

  \caption{
  Human frame generation and selection. The image model inserts one person into the original ScanNet++ view according to the interaction instruction while preserving the scene, camera viewpoint, and target object. When multiple candidates are generated, a vision-language evaluator selects the frame that best preserves the scene and expresses the required body--scene contacts. The selected frame is used as an intermediate visual hypothesis for GVHMR initialization and contact estimation.
  }
  \label{fig:human-frame-generation}
\end{figure}

\section{Agentic Diffusion-Based Contact Estimation}
\label{sec:method-contact-estimation}

The third module estimates the scene-side contact regions required by the interaction. This stage is the main bridge between open-vocabulary visual reasoning and metric 3D optimization. Given the selected human frame $\hat{I}$, the original scene image $I$, and the Scene Interaction Graph $G$, the goal is to obtain a set of body-part-specific contact masks on the original scene image. These masks are later projected to the reconstructed scene mesh and used as 3D contact targets.

The key difficulty is that the image model can reason about contacts from the generated human frame, but its edited outputs are not geometrically reliable by themselves. When asked to draw contact overlays, the model may slightly change the target object, modify the background, alter image texture, or shift scene boundaries. The method therefore does not use the generated overlay image as the final annotated image. Instead, it extracts only the solid-color contact regions from the generated output and pastes those regions back onto the unchanged original scene image. The resulting image is called the contact composite. This ensures that the verified masks are always defined in the coordinate system of the calibrated input image $I$.

\subsection{Reference, Canvas, and Contact-Overlay Images}
\label{sec:method-contact-reference-canvas}

For contact estimation, the selected human frame $\hat{I}$ serves as the reference image. It contains the generated person and therefore shows the approximate pose, body laterality, and visible contact configuration. The original ScanNet++ image $I$ serves as the canvas image. It contains the scene without the generated human and preserves the original target-object appearance, location, and boundaries.

For every object-contact edge $e = (h_e, s_e)$ in the SIG, the model must identify the localized scene surface contacted by body part $h_e$. If the instruction is ``a person opening a washing machine,'' the corresponding hand contact should be localized near the functional door or handle region. If the instruction specifies a treadmill interaction with both hands on the side frames and one foot on the belt, the contact regions should be localized on the treadmill rails and belt rather than arbitrary treadmill surfaces.

Each required contact edge is assigned a distinct bright color. Let $c_e$ denote the color assigned to edge $e$. At iteration $r$, the image model receives the reference image $\hat{I}$, the canvas image $I$, the target-object description, the current correction instruction if one exists, and the color mapping from contact edges to colors. It produces an overlay candidate
\begin{equation}
\label{eq:contact-overlay-candidate}
O^{(r)}
=
\Phi_{\mathrm{con}}
(
\hat{I},
I,
G,
\{c_e\},
a^{(r-1)}
),
\end{equation}
where $\Phi_{\mathrm{con}}$ denotes the contact-overlay generation model and $a^{(r-1)}$ is the correction instruction from the previous iteration. For the first iteration, no correction instruction is used. The full contact-overlay prompt is provided in \cref{app:system-prompt-contact-overlay}.

The intended output $O^{(r)}$ is an edited copy of the canvas image with solid-color masks drawn only on the localized contact regions of the target object. However, because the generated image may also contain unintended changes to the scene, $O^{(r)}$ is treated only as a temporary source from which color masks are extracted.

\subsection{Overlay Extraction and Composite Construction}
\label{sec:method-contact-composite}

After the image model produces the overlay candidate $O^{(r)}$, the method extracts the pixels corresponding to each assigned contact color. For each object-contact edge $e$, the color extraction step produces a binary mask
\begin{equation}
\label{eq:color-extraction}
Z_e^{(r)}
=
\operatorname{ExtractColor}
(
O^{(r)},
c_e
).
\end{equation}
The extracted mask $Z_e^{(r)}$ contains only pixels whose color matches the assigned body-part color $c_e$, within a small tolerance. This step discards all other content in the generated overlay image, including any accidental modifications to scene texture, object shape, lighting, or background.

The extracted masks are then placed on top of the original unedited scene image $I$. This produces the contact composite
\begin{equation}
\label{eq:contact-composite}
C^{(r)}
=
\operatorname{Composite}
(
I,
\{Z_e^{(r)}\}
).
\end{equation}
The composite $C^{(r)}$ has the same camera coordinate system, resolution, and scene appearance as the original image $I$. The only difference is the addition of solid-color contact overlays. This design is important for projection: the final masks must correspond to pixels in the calibrated original image, not to a potentially altered image produced by the generative model.

The composite also separates two roles that would otherwise be entangled. The image generation model is used to infer and draw candidate contact regions, while the original scene image remains the authoritative geometric canvas. Thus, even if the overlay candidate $O^{(r)}$ contains minor scene drift, the downstream verifier and projector only operate on the clean composite $C^{(r)}$.

\subsection{Agentic Contact Verification and Correction}
\label{sec:method-agentic-contact-verification}

The extracted contact overlays are verified using a closed-loop vision-language evaluation stage. The evaluator receives three images: the selected human frame $\hat{I}$, the original scene image $I$, and the contact composite $C^{(r)}$. The selected human frame shows the intended human-object interaction, the original scene image provides the unmodified target-object geometry, and the composite shows the current proposed contact masks on that original scene.

The evaluator checks whether the colored overlays in $C^{(r)}$ correctly mark the scene regions contacted by the generated human in $\hat{I}$. It verifies that all required contacts are present, that no extra contacts are added, that each color corresponds to the correct body part, that each mask lies on the correct object surface, and that each mask has an appropriate size and location. The evaluation prompt is provided in \cref{app:system-prompt-contact-eval}.

The evaluator returns either an acceptance decision or a correction instruction. This can be written as
\begin{equation}
\label{eq:vlm-contact-eval}
(d^{(r)}, a^{(r)})
=
\Phi_{\mathrm{eval}}
(
\hat{I},
I,
C^{(r)},
G,
\{c_e\}
),
\end{equation}
where $d^{(r)}$ is the decision and $a^{(r)}$ is a natural-language correction instruction. If the masks are correct, $d^{(r)}$ is accepted and the loop terminates. If the masks are rejected, the correction instruction is passed back to the contact-overlay generation model for the next iteration. For example, the instruction may ask the model to move a hand mask to the correct handle region, shrink an overly large contact footprint, add a missing contact, remove an extra mask, or correct a body-part color.

This agentic loop is repeated for a limited number of iterations:
\begin{equation}
\label{eq:agentic-contact-loop}
O^{(r)}
\rightarrow
\{Z_e^{(r)}\}
\rightarrow
C^{(r)}
\rightarrow
(d^{(r)}, a^{(r)})
\rightarrow
O^{(r+1)}.
\end{equation}
The first accepted composite is used as the final contact annotation. If no candidate is accepted within the maximum number of attempts, the best available composite is retained only when the contact locations are still usable for optimization; otherwise the interaction is marked as a failure case.

Let $r^\star$ denote the selected iteration. The final object-contact masks are
\begin{equation}
\label{eq:final-object-contact-masks}
M_e^{2D}
=
Z_e^{(r^\star)}
\qquad
\text{for all}
\qquad
e \in \mathcal{E}_{obj},
\end{equation}
where $\mathcal{E}_{obj}$ is the subset of SIG edges whose scene node is the target object. Small isolated components may be removed after extraction, and masks can be clipped to the visible target-object region when necessary.

\subsection{Floor Support Masks}
\label{sec:method-floor-support-masks}

Floor contacts are handled separately from object contacts. For support edges such as left-foot--floor or right-foot--floor, the exact pixel location of the foot in the original image is less important than obtaining a geometrically valid support surface. The method therefore uses SAM3 to segment visible floor regions in the original scene image $I$. The resulting floor mask is projected to the reconstructed scene mesh and used as the scene-side support target for foot contacts.

This separation avoids asking the contact-overlay model to hallucinate precise foot footprints on an empty floor region in the original image. For standing, walking, or balancing interactions, the projected floor region provides a stable support surface, while the SMPL-X optimization determines the final precise foot placement subject to contact, penetration, and pose priors.

Let $\mathcal{E}_{floor}$ denote the subset of SIG edges whose scene node is the floor. For each floor-support edge $e$, the corresponding 2D support mask is denoted by
\begin{equation}
\label{eq:floor-mask-2d}
M_e^{2D}
=
M_{floor}^{2D}
\qquad
\text{for}
\qquad
e \in \mathcal{E}_{floor}.
\end{equation}
The same floor mask can be reused for multiple foot-support edges, since the distinction between left and right foot is enforced on the body side during optimization.

\subsection{Projection to 3D Contact Targets}
\label{sec:method-contact-projection}

The verified 2D contact masks are projected to the reconstructed scene mesh using the calibrated camera. For a 3D scene point $x$ in the ScanNet++ world coordinate frame, its image projection is
\begin{equation}
\label{eq:scene-point-projection}
\tilde{u}
=
K
(
R x + t
),
\end{equation}
followed by perspective division to obtain the pixel coordinate $u = (u_x,u_y)$. A scene vertex or sampled scene point is assigned to contact edge $e$ if its projected pixel lies inside the final mask $M_e^{2D}$ and passes visibility filtering for the selected camera view.

This projection converts each image-space mask into a metric scene-side contact target:
\begin{equation}
\label{eq:projected-contact-target}
P_e^s
=
\{x \in M_s \mid \pi(x;K,R,t) \in M_e^{2D}\}.
\end{equation}
Here, $\pi(x;K,R,t)$ denotes the calibrated projection from the scene coordinate frame to the image plane. In practice, $P_e^s$ is represented as a set of scene vertices or sampled surface points on the reconstructed mesh. These points remain fixed during optimization and serve as the target regions that the corresponding SMPL-X contact vertices should approach.

The final output of the contact-estimation module is the set of projected contact targets
\begin{equation}
\label{eq:projected-contact-target-set}
\mathcal{P}^{s}
=
\{P_e^s \mid e \in \mathcal{E}\}.
\end{equation}
Each element of this set is linked to one SIG edge and therefore to one body-side SMPL-X contact region. The next modules use these targets to initialize and optimize the human body in the reconstructed scene.

\begin{figure}[t]
  \centering
  % TODO: Replace with real panels from one example.
  % Recommended panels:
  % (a) original ScanNet++ image I
  % (b) selected generated human frame \hat{I}
  % (c) generated contact-overlay candidate O^{(r)}
  % (d) extracted masks pasted onto original image: composite C^{(r)}
  % (e) VLM correction / accepted composite
  % (f) projected 3D contact targets on the scene mesh

  \fbox{
    \begin{minipage}[c][0.28\textheight][c]{0.94\textwidth}
      \centering
      \textbf{Agentic contact estimation placeholder}\\[0.8em]
      Human frame $\hat{I}$
      \quad $+$ \quad
      original image $I$
      \quad $\rightarrow$ \quad
      overlay candidate $O^{(r)}$
      \quad $\rightarrow$ \quad
      extracted masks $Z_e^{(r)}$
      \quad $\rightarrow$ \quad
      composite $C^{(r)}$\\[0.8em]
      VLM verification
      \quad $\rightarrow$ \quad
      correction loop or accepted masks
      \quad $\rightarrow$ \quad
      3D contact targets $P_e^s$
    \end{minipage}
  }

  \caption{
  Agentic diffusion-based contact estimation. The image model proposes colored contact overlays using the selected human frame as reference and the original scene image as canvas. Since the generated overlay image may still alter the underlying scene appearance or geometry, only the colored mask pixels are extracted and recomposited onto the untouched original image. A vision-language evaluator checks the composite against the human frame and original scene image, then either accepts the masks or returns a correction instruction for another generation attempt. Accepted masks are projected to the reconstructed scene mesh to obtain body-part-specific 3D contact targets.
  }
  \label{fig:agentic-contact-estimation}
\end{figure}

\section{SMPL-X Pose Initialization}
\label{sec:method-pose-init}

The fourth module initializes the human body from the selected generated frame. The contact-estimation stage provides scene-side contact targets, but these targets alone do not determine a complete articulated body. The method therefore uses GVHMR to obtain an initial SMPL-X estimate from the generated human image $\hat{I}$. This initialization provides a plausible full-body pose, body orientation, translation, and shape before scene-aware refinement.

GVHMR is designed for monocular human reconstruction from image or video input. Since the proposed setting uses a single static generated frame, the selected image $\hat{I}$ is converted into a short static clip by repeating the same frame:
\begin{equation}
\label{eq:static-human-clip}
\mathcal{V}_{\hat{I}}
=
(
\hat{I},
\hat{I},
\ldots,
\hat{I}
).
\end{equation}
This allows the frame to be processed by the GVHMR pipeline while preserving the static nature of the target interaction. The first reconstructed frame is used as the initialization for the downstream optimization.

Applying GVHMR to the static clip produces an initial SMPL-X estimate in the camera coordinate frame:
\begin{equation}
\label{eq:gvhmr-camera-init}
\theta^{cam}_0
=
(
\mathbf{T}^{cam}_0,
\mathbf{r}^{cam}_0,
\mathbf{p}_0,
\boldsymbol{\beta}_0
).
\end{equation}
Here, $\mathbf{T}^{cam}_0$ is the initial global body translation in camera coordinates, $\mathbf{r}^{cam}_0$ is the initial global body orientation, $\mathbf{p}_0$ is the articulated body pose, and $\boldsymbol{\beta}_0$ is the estimated SMPL-X body shape. The shape estimate is used as a fixed body identity during optimization.

The final optimization is performed in the ScanNet++ scene coordinate frame, so the GVHMR initialization is transformed from camera coordinates to world coordinates using the selected image camera pose. With the projection convention
\begin{equation}
\label{eq:camera-transform-convention}
x^{cam}
=
R x^{world} + t,
\end{equation}
the inverse transform maps camera-frame points into the ScanNet++ world frame. The initial translation is converted as
\begin{equation}
\label{eq:init-translation-world}
\mathbf{T}_0
=
R^{\top}
(
\mathbf{T}^{cam}_0 - t
).
\end{equation}
The global body orientation is converted by applying the inverse camera rotation to the camera-frame body orientation. If $\mathbf{R}^{cam}_0$ denotes the rotation matrix corresponding to $\mathbf{r}^{cam}_0$, the world-frame orientation is
\begin{equation}
\label{eq:init-orientation-world}
\mathbf{R}_0
=
R^{\top}
\mathbf{R}^{cam}_0.
\end{equation}
The corresponding axis-angle representation is denoted by $\mathbf{r}_0$.

The resulting initialization in the scene coordinate frame is
\begin{equation}
\label{eq:gvhmr-world-init}
\theta_0
=
(
\mathbf{T}_0,
\mathbf{r}_0,
\mathbf{p}_0,
\boldsymbol{\beta}_0
).
\end{equation}
This initialization is visually plausible because it is estimated from the generated human frame, but it is not guaranteed to satisfy the projected contact targets or avoid collisions with the reconstructed scene mesh. In particular, the generated image may contain small inconsistencies in scale, depth, contact location, or body placement. The role of the next stage is therefore to refine $\mathbf{T}_0$, $\mathbf{r}_0$, and $\mathbf{p}_0$ using the scene geometry and the contact targets obtained from the previous module.

During optimization, the shape parameters $\boldsymbol{\beta}_0$ are kept fixed. This avoids using contact refinement to compensate for geometric errors by changing body identity. The optimized variables are limited to global translation, global orientation, and articulated body pose, while the reconstructed scene mesh and projected contact targets remain fixed.

\section{Scene-Aware SMPL-X Optimization}
\label{sec:method-optimization}

The final module refines the initialized SMPL-X body using the reconstructed scene geometry and the projected contact targets. The GVHMR estimate provides a plausible human pose from the generated image, but it is not guaranteed to satisfy the contact requirements in 3D. In particular, the body may be slightly offset from the target object, may contact an incorrect surface, or may intersect the reconstructed scene. The optimization stage corrects these errors while preserving the plausibility of the initial body pose.

\subsection{Optimized Variables}
\label{sec:method-optimized-variables}

The optimization starts from the world-frame initialization
\begin{equation}
\label{eq:optimization-init}
\theta_0
=
(
\mathbf{T}_0,
\mathbf{r}_0,
\mathbf{p}_0,
\boldsymbol{\beta}_0
).
\end{equation}
Only the global translation $\mathbf{T}$, global orientation $\mathbf{r}$, and articulated body pose $\mathbf{p}$ are optimized. The body shape $\boldsymbol{\beta}_0$ is kept fixed throughout the optimization. The optimized body parameters are therefore
\begin{equation}
\label{eq:optimized-theta}
\theta
=
(
\mathbf{T},
\mathbf{r},
\mathbf{p},
\boldsymbol{\beta}_0
).
\end{equation}
Keeping the shape fixed prevents the optimizer from satisfying contact constraints by changing the body identity or body proportions. The optimization is instead forced to explain the interaction through pose and placement.

For any current parameter state $\theta$, the SMPL-X model produces a posed human mesh
\begin{equation}
\label{eq:optimized-human-mesh}
M_h(\theta)
=
(
V_h(\theta),
F_h
).
\end{equation}
The scene mesh $M_s$ and the projected contact targets $P_e^s$ remain fixed. Thus, the optimization changes only the human body, not the reconstructed scene or the estimated scene-side contact regions.

\subsection{Contact Objective}
\label{sec:method-contact-objective}

The contact objective attracts the specified SMPL-X body regions toward the corresponding projected scene contact targets. For each SIG edge $e$, the body-side contact vertices are given by the fixed vertex subset $B_e$, and the corresponding scene-side target is $P_e^s$. The body-side contact points under the current pose are
\begin{equation}
\label{eq:body-contact-current}
P_e^h(\theta)
=
\{V_h(\theta)_i \mid i \in B_e\}.
\end{equation}

The contact loss is defined as a nearest-neighbor distance from the selected body contact vertices to the projected scene contact targets:
\begin{equation}
\label{eq:contact-loss}
\mathcal{L}_{\mathrm{contact}}(\theta)
=
\sum_{e \in \mathcal{E}}
\frac{1}{|B_e|}
\sum_{i \in B_e}
\operatorname*{min}_{q \in P_e^s}
\left\|
V_h(\theta)_i - q
\right\|_2^2.
\end{equation}
This loss is edge-specific: hand vertices are pulled toward hand contact targets, foot vertices toward floor or object support targets, and so on. The SIG therefore determines which body region should be attracted to which scene region, while the projected masks determine the metric 3D location of the target.

The one-sided nearest-neighbor formulation is useful because the scene target can contain multiple vertices or sampled surface points. The body contact region is encouraged to approach the nearest compatible part of the target region without requiring a fixed point-to-point correspondence between SMPL-X vertices and scene vertices.

\subsection{Penetration Objective}
\label{sec:method-penetration-objective}

Contact attraction alone can produce physically invalid solutions if the body moves into the scene geometry. To discourage intersections, the optimization includes a penetration objective based on a signed distance representation of the SMPL-X body. The method uses VolumetricSMPL to query whether sampled scene points lie inside the current posed body.

Let $\mathcal{Q}_s$ denote a set of scene points sampled from the reconstructed mesh near the current body. Let $\phi_{\theta}(x)$ be the signed distance from a scene point $x$ to the posed SMPL-X body, with negative values indicating that $x$ lies inside the body volume. The penetration loss is
\begin{equation}
\label{eq:penetration-loss}
\mathcal{L}_{\mathrm{penetration}}(\theta)
=
\frac{1}{|\mathcal{Q}_s|}
\sum_{x \in \mathcal{Q}_s}
\left[
\operatorname*{max}
(
0,
-\phi_{\theta}(x)
)
\right]^2.
\end{equation}
This term penalizes scene points that are inside the human body and assigns zero cost to points outside the body. It therefore discourages severe body-scene interpenetration while still allowing nearby contact.

\subsection{Pose and Orientation Priors}
\label{sec:method-pose-orientation-priors}

The contact and penetration losses are not sufficient to guarantee a natural human pose. Since the initialization from GVHMR is visually plausible, the optimization includes priors that keep the refined body close to the initial estimate. The pose prior penalizes large deviations from the initial articulated pose:
\begin{equation}
\label{eq:pose-prior-loss}
\mathcal{L}_{\mathrm{pose}}(\theta)
=
\left\|
\mathbf{p}
-
\mathbf{p}_0
\right\|_2^2.
\end{equation}

The orientation prior penalizes large changes in the global body orientation. Let $\mathbf{R}(\mathbf{r})$ and $\mathbf{R}(\mathbf{r}_0)$ be the rotation matrices corresponding to the current and initial global orientations. The orientation prior is defined using the geodesic distance on rotations:
\begin{equation}
\label{eq:orientation-prior-loss}
\mathcal{L}_{\mathrm{orient}}(\theta)
=
\left\|
\operatorname{Log}
\left(
\mathbf{R}(\mathbf{r}_0)^{\top}
\mathbf{R}(\mathbf{r})
\right)
\right\|_2^2.
\end{equation}
Here, $\operatorname{Log}$ maps a rotation matrix to its axis-angle vector. This term prevents the optimizer from satisfying contacts through unrealistic global rotations that are inconsistent with the generated human frame.

\subsection{Total Objective}
\label{sec:method-total-objective}

The final optimization objective combines contact attraction, penetration avoidance, and the initialization priors:
\begin{equation}
\label{eq:optimization-total-loss}
\mathcal{L}(\theta)
=
\lambda_{\mathrm{contact}}
\mathcal{L}_{\mathrm{contact}}(\theta)
+
\lambda_{\mathrm{penetration}}
\mathcal{L}_{\mathrm{penetration}}(\theta)
+
\lambda_{\mathrm{pose}}
\mathcal{L}_{\mathrm{pose}}(\theta)
+
\lambda_{\mathrm{orient}}
\mathcal{L}_{\mathrm{orient}}(\theta).
\end{equation}
The optimized parameters are obtained by
\begin{equation}
\label{eq:optimization-argmin}
\theta^{\ast}
=
\operatorname*{arg\,min}_{\mathbf{T},\,\mathbf{r},\,\mathbf{p}}
\mathcal{L}(\theta),
\qquad
\boldsymbol{\beta} = \boldsymbol{\beta}_0.
\end{equation}

The contact term enforces the functional body--scene relationships inferred by the previous modules. The penetration term prevents the body from achieving contact by moving through scene geometry. The pose and orientation priors preserve the plausibility of the GVHMR initialization. Together, these terms refine the initial SMPL-X estimate into a geometrically grounded static interaction state.

\subsection{Optimization Procedure}
\label{sec:method-optimization-procedure}

The optimization is initialized with $\theta_0$ from GVHMR. At each iteration, the SMPL-X mesh is recomputed, body contact vertices are matched to their corresponding scene contact targets, penetration is evaluated against nearby scene samples, and gradients are backpropagated through the differentiable SMPL-X model. The scene mesh, projected contact targets, and body shape remain fixed.

The optimization returns the final parameters
\begin{equation}
\label{eq:final-optimized-body}
\theta^{\ast}
=
(
\mathbf{T}^{\ast},
\mathbf{r}^{\ast},
\mathbf{p}^{\ast},
\boldsymbol{\beta}_0
).
\end{equation}
The final output mesh is
\begin{equation}
\label{eq:final-human-mesh}
M_h^{\ast}
=
M_h(\theta^{\ast}),
\end{equation}
which is already expressed in the ScanNet++ scene coordinate frame. This mesh is the final synthesized static human--scene interaction result.

The optimization is carried out using Adam for a fixed budget of 2000 iterations. This fixed optimization budget is used across all interactions.